% Options for packages loaded elsewhere
\PassOptionsToPackage{unicode}{hyperref}
\PassOptionsToPackage{hyphens}{url}
\PassOptionsToPackage{dvipsnames,svgnames,x11names}{xcolor}
%
\documentclass[
  11pt,
  a4paper,
]{ltjsarticle}
\usepackage{amsmath,amssymb}
\usepackage{iftex}
\ifPDFTeX
  \usepackage[T1]{fontenc}
  \usepackage[utf8]{inputenc}
  \usepackage{textcomp} % provide euro and other symbols
\else % if luatex or xetex
  \usepackage{unicode-math} % this also loads fontspec
  \defaultfontfeatures{Scale=MatchLowercase}
  \defaultfontfeatures[\rmfamily]{Ligatures=TeX,Scale=1}
\fi
\usepackage{lmodern}
\ifPDFTeX\else
  % xetex/luatex font selection
\fi
% Use upquote if available, for straight quotes in verbatim environments
\IfFileExists{upquote.sty}{\usepackage{upquote}}{}
\IfFileExists{microtype.sty}{% use microtype if available
  \usepackage[]{microtype}
  \UseMicrotypeSet[protrusion]{basicmath} % disable protrusion for tt fonts
}{}
\makeatletter
\@ifundefined{KOMAClassName}{% if non-KOMA class
  \IfFileExists{parskip.sty}{%
    \usepackage{parskip}
  }{% else
    \setlength{\parindent}{0pt}
    \setlength{\parskip}{6pt plus 2pt minus 1pt}}
}{% if KOMA class
  \KOMAoptions{parskip=half}}
\makeatother
\usepackage{xcolor}
\usepackage[margin=24mm]{geometry}
\setlength{\emergencystretch}{3em} % prevent overfull lines
\providecommand{\tightlist}{%
  \setlength{\itemsep}{0pt}\setlength{\parskip}{0pt}}
\setcounter{secnumdepth}{5}
\ifLuaTeX
  \usepackage{selnolig}  % disable illegal ligatures
\fi
\IfFileExists{bookmark.sty}{\usepackage{bookmark}}{\usepackage{hyperref}}
\IfFileExists{xurl.sty}{\usepackage{xurl}}{} % add URL line breaks if available
\urlstyle{same}
\hypersetup{
  colorlinks=true,
  linkcolor={blue},
  filecolor={Maroon},
  citecolor={Blue},
  urlcolor={blue},
  pdfcreator={LaTeX via pandoc}}

\author{}
\date{}

\begin{document}

\hypertarget{ux6b21ux5143}{%
\section{次元}\label{ux6b21ux5143}}

有限次元ベクトル空間 \texttt{V}
の次元は、基底に含まれるベクトルの本数です。

\[
\dim V=n
\]

とは、\texttt{V} を表現するために独立な方向が \texttt{n}
個必要だということです。

\hypertarget{ux6210ux5206ux6570ux3068ux672cux8ceaux7684ux306aux6b21ux5143ux306fux9055ux3046}{%
\subsection{「成分数」と「本質的な次元」は違う}\label{ux6210ux5206ux6570ux3068ux672cux8ceaux7684ux306aux6b21ux5143ux306fux9055ux3046}}

1000個の成分を持つデータでも、そのデータが実際には10個の独立な方向の組合せでほぼ表せるなら、本質的な自由度は10程度かもしれません。

この考え方が

\begin{itemize}
\tightlist
\item
  ランク
\item
  低ランク近似
\item
  PCA
\item
  manifold 的なデータ理解
\end{itemize}

につながります。

\hypertarget{ux6b21ux5143ux306fux57faux5e95ux306bux3088ux3089ux306aux3044}{%
\subsection{次元は基底によらない}\label{ux6b21ux5143ux306fux57faux5e95ux306bux3088ux3089ux306aux3044}}

異なる基底を選んでも基底ベクトルの本数は同じです。したがって次元は座標系ではなく、空間そのものに固有の量です。

\end{document}
